\documentclass[11pt]{article}
\usepackage[margin=1in]{geometry}
\usepackage{amsmath,amssymb,amsthm}
\usepackage[hidelinks]{hyperref}
\usepackage{graphicx}

\newtheorem{theorem}{Theorem}
\newtheorem{lemma}{Lemma}
\newtheorem{remark}{Remark}

\newcommand{\Kinf}{K_{\infty,1}}
\newcommand{\phic}{\varphi^c}
\newcommand{\phiz}{\varphi^0}

\title{A diagonal-couple counterexample to pairwise necessity of \(\Kinf\)}
\author{Run fa\_banach\_001, lane 19}
\date{2026-06-20}

\begin{document}
\maketitle

\section*{Source problem}

The source paper is Yves Raynaud and Pedro Tradacete, \emph{Calderon-Lozanovskii interpolation on quasi-Banach lattices}, arXiv:1611.09079. Corollaries 5.1 and 5.6 prove interpolation-functor conclusions for the methods \(\phic\) and \(\phiz\) when the target endpoint quasi-Banach lattices have the \(\Kinf\) property. Remark 5.8 asks:
\[
\text{``We do not know whether the \(\Kinf\) property in Corollaries 5.1 and 5.6 is actually necessary.''}
\]

\begin{center}
\includegraphics[width=0.92\linewidth]{figures/open_problem_crop.png}
\end{center}

\section*{Idea of the proof}

The source paper already supplies a quasi-Banach lattice \(E\) without \(\Kinf\) in Example 4.3. Put that lattice on both endpoints. Then the Calderon-Lozanovskii construction has no room to create a genuinely new target space: for every normalized \(\varphi\), the diagonal space \(\varphi(E,E)\) is just \(E\) with an equivalent quasi-norm. The Gagliardo completion \(\phic(E,E)\) and the intersection closure \(\phiz(E,E)\) are again \(E\). Thus any operator bounded into both endpoints is simply bounded into \(E\), and the interpolation conclusions become automatic.

The counterexample is therefore intentionally degenerate. It disproves pairwise necessity for the corollary conclusions, but it does not classify all target couples and does not rule out a stronger necessity theorem after excluding diagonal collapse or imposing a non-degeneracy condition.

\section*{Definitions used}

A quasi-Banach lattice \(E\) has the \(\Kinf\) property if there is \(C>0\) such that for every finite family \(x_1,\ldots,x_n\in E\),
\[
\left\|\bigvee_{i=1}^n |x_i|\right\|_E
\leq C \max_{|a_i|\leq 1}\left\|\sum_{i=1}^n a_i x_i\right\|_E.
\]
This is the equivalent intrinsic form recorded in Proposition 4.2 of the source paper. Example 4.3 constructs a quasi-Banach lattice \(E\) which fails this property.

We use the source paper's usual normalization \(\varphi(1,1)=1\) for functions \(\varphi\in\mathcal P\). Multiplying \(\varphi\) by a positive scalar only changes the resulting quasi-norms by an equivalent factor, so this normalization does not affect the necessity issue.

\begin{lemma}[Diagonal Calderon-Lozanovskii collapse]
Let \(E\) be a quasi-Banach lattice and let \(\varphi\in\mathcal P\) be normalized by \(\varphi(1,1)=1\). Then
\[
\varphi(E,E)\simeq E,\qquad \phic(E,E)\simeq E,\qquad \phiz(E,E)\simeq E.
\]
\end{lemma}

\begin{proof}
Let \(A_E\geq 1\) be a quasi-triangle constant for \(E\). Since \(\varphi\) is increasing in both variables and homogeneous, for scalars \(s,t\geq 0\) one has
\[
\varphi(s,t)\leq \max(s,t).
\]
Indeed, if \(s,t\leq m\), then \(\varphi(s,t)\leq\varphi(m,m)=m\varphi(1,1)=m\).

Suppose \(\|x\|_{\varphi(E,E)}<\lambda\). Then there are \(u_0,u_1\in E_+\) with \(\|u_0\|_E,\|u_1\|_E\leq 1\) and
\[
|x|\leq \lambda\,\varphi(u_0,u_1)\leq \lambda\,(u_0\vee u_1).
\]
By lattice monotonicity and the quasi-triangle inequality,
\[
\|x\|_E\leq \lambda\|u_0\vee u_1\|_E
\leq \lambda\|u_0+u_1\|_E
\leq 2A_E\lambda.
\]
Hence the identity map \(\varphi(E,E)\to E\) is continuous.

Conversely, if \(x\neq 0\), put \(u=|x|/\|x\|_E\). Then \(\|u\|_E=1\) and
\[
|x|=\|x\|_E\,\varphi(u,u),
\]
so \(\|x\|_{\varphi(E,E)}\leq\|x\|_E\). Thus \(\varphi(E,E)\) and \(E\) are the same lattice with equivalent quasi-norms.

For the diagonal couple, the ambient sum \(E+E\) is \(E\). The Gagliardo completion \(\phic(E,E)\) is obtained by closing the \(\varphi(E,E)\)-unit ball inside \(E\). Since the just-proved equivalence sandwiches that unit ball between two \(E\)-balls, the resulting quasi-norm is again equivalent to the norm of \(E\). Finally, \(E\cap E=E\), so the closure defining \(\phiz(E,E)\) is also \(E\) with the same equivalent topology.
\end{proof}

\begin{theorem}
The \(\Kinf\) property is not necessary for the conclusions of Corollaries 5.1 and 5.6 of the source paper in the literal pairwise target-couple sense.
\end{theorem}

\begin{proof}
Let \(E\) be the quasi-Banach lattice without \(\Kinf\) constructed in Example 4.3 of the source paper, and consider the target couple \((Y_0,Y_1)=(E,E)\). We verify the conclusions of the two corollaries for this target couple.

Let \((X_0,X_1)\) be any compatible couple of quasi-Banach lattices and let
\[
T:(X_0,X_1)\to(E,E)
\]
be endpoint bounded. Write \(M_j=\|T:X_j\to E\|\). If \(x=x_0+x_1\in X_0+X_1\), then
\[
\|Tx\|_E\leq A_E\bigl(\|Tx_0\|_E+\|Tx_1\|_E\bigr)
\leq A_E\max(M_0,M_1)\bigl(\|x_0\|_{X_0}+\|x_1\|_{X_1}\bigr).
\]
Taking the infimum over all decompositions \(x=x_0+x_1\) shows that \(T:X_0+X_1\to E\) is bounded.

By the definition of the Gagliardo completion used in the source paper, \(\phic(X_0,X_1)\) embeds continuously into \(X_0+X_1\). Hence
\[
T:\phic(X_0,X_1)\to E
\]
is bounded. The lemma identifies \(E\simeq\phic(E,E)\), giving the conclusion of Corollary 5.1 for this target couple. Since the source paper records bounded inclusions
\[
\phiz(X_0,X_1)\subset \varphi(X_0,X_1)\subset \phic(X_0,X_1),
\]
the same bounded operator restricts to
\[
T:\phiz(X_0,X_1)\to E\simeq\phiz(E,E),
\]
which is the conclusion of Corollary 5.6 for this target couple.

Nevertheless, \(Y_0=Y_1=E\) fails \(\Kinf\). Thus \(\Kinf\) is not necessary for the stated functorial conclusions when necessity is read pairwise for target couples.
\end{proof}

\section*{Verification notes and limitations}

The proof uses only the diagonal-couple collapse and the non-\(\Kinf\) example from the source paper. It does not claim that \(\Kinf\) is unnecessary for all target couples, and it does not answer a stronger classification problem asking which non-diagonal target couples satisfy the \(\phic\) or \(\phiz\) interpolation conclusions.

Cheap run indexes were searched for the arXiv id and core keywords. A bounded web search on 2026-06-20 for exact \(K_{\infty,1}\)/Calderon-Lozanovskii necessity phrases only surfaced the source paper. No separate later solution or duplicate packet was found.

\section*{References}

\begin{enumerate}
\item Yves Raynaud and Pedro Tradacete, \emph{Calderon-Lozanovskii interpolation on quasi-Banach lattices}, arXiv:1611.09079.
\end{enumerate}

\end{document}
